\documentclass{article}
\usepackage{amsmath}
\usepackage{listings}
\usepackage{hyperref}
\hypersetup{
  colorlinks=true,
  linkcolor=blue,
  filecolor=magenta,
  urlcolor=cyan,
  pdftitle={PMFRG For Dummies},
  pdfpagemode=FullScreen,
}

\urlstyle{same}

\title{PMFRG for dummies: Useful formulas for the implementation of the PMFRG code}
\author{michele.mesiti }
\date{July 2026}

\begin{document}

\maketitle

\section{What's this ``document''?}

In this document all the necessary equations
for the implementation of the PMFRG method
are written in the most explicit way possible.

These notes are work in progress,
probably will never be really finished.

A more modern approach would suggest to write this up
as part of the documentation of the code,
to be managed with Documenter.jl,
but at the moment LaTeX might still fit better.
This is not set in stone.
Maybe it would be better to use typst?
In any case hopefully Pandoc will prove once again its worth.

At the moment this document can be rendered with latex
with the shell command

\begin{lstlisting}
$ pdflatex notex.tex
\end{lstlisting}

or

\begin{lstlisting}
$ make
\end{lstlisting}



\section{Preliminary notes}

\subsection{Notation: Multiindices}
\label{multiindex}

In various context different kinds of multiindices are used.
In the general review [\ref{PFFRGReview}], just after Eq. (16)
multiindices are introduced as $x = (i,i\omega,\mu)$,
where:

\begin{itemize}
\item $i$ is the lattice site index
\item $\omega$ is the Matsubara frequency
\item $\mu$ is the ``spin index'' (in the PFFRG case), which
  generalizes to the ``flavour index''.
\end{itemize}

In what follows, thought, we want to be rigorous,
and rigor is best communicate by being explicit,
so we will use the meaning of the multiindices
described above
to translate equations to the dumbest, most explicit form.

\subsection{Propagators}

In what follows:
\begin{itemize}
\item $a$, $b$ and latin lower-case letters represent Majorana flavours;
\item $x$ is the site index ($G$ can depend only on one site, because of the local $Z_2$ gauge symmetry - see \ref{PFFRGReview});
\item $T$ is the temperature, acting in the $T$-flow case as a flow parameter;
\item $\Lambda$ the cutoff scale (only in $\Lambda$-flow).
\end{itemize}

Note: In the implementation, $\Sigma$ does not have an explicit dependency on the flow parameter
(the \textsf{iSigma\_} functions just retrieve values from a data structure with
minimal manipulation).
For this reason,
in for $iG$ we use the flow parameter as a function argument,
while in the case of $\Sigma$ we merely add it as a superscript.


\subsubsection{Matrix algebra for dummies: Derivative of $M^{-1}$}
Notice that $G$ is represented as the inverse of a matrix.
The derivative of the inverse of a matrix $M$ can be computed by noticing that
\begin{equation}
  \partial_\Lambda\left(M^{-1}M\right) = 0
\end{equation}
since $M^{-1}M = M M^{-1} = I$ (assuming $M$ is invertible, of course).
This means
\begin{equation}
  \left(\partial_\Lambda M^{-1}\right)  M + M^{-1}\left(\partial_\Lambda M\right) = 0 .
\end{equation}
So
\begin{equation}
  \label{derivinv}
 \frac{\partial M}{\partial \Lambda} = - M \frac{\partial M^{-1}}{\partial \Lambda}  M,
\end{equation}
which reduces to the simple case
\begin{equation}
 \frac{\partial M}{\partial \Lambda} = - M^2 \frac{\partial M^{-1}}{\partial \Lambda}
\end{equation}
if the derivative of $M^{-1}$ commutes with $M$.


\subsubsection{$T$ flow}
The propagator depends on $\Sigma$:
\begin{equation}
  \label{g-tflow-general}
  i{(G(\omega,T,x))}^{-1}_{ab} = \left(\sqrt{T} \omega\right) \delta_{ab} +  i\Sigma^T_{ab}(x,\omega)
\end{equation}


In the isotropic (Heisenberg) or diagonal (``XYZ'') case
the inverse is trivial, but not in the most general anisotropic case,
where some matrix inversion function needs to be used.

The expression \ref{g-tflow-general} comes from \href{https://github.com/NoahHassan/pmfrg\_xyz/blob/1824a2bf7811b5276eee6d8d7f40087b5d8a4c15/src/PMFRG_general.jl\#L146-L159}{Noah's general PMFRG}
code, but there should be a more official source.

From the $T$-$\Lambda$ flow [\ref{PMFRGXYZTLambda}] (see line 1599) we get the same expression.

\textbf{TODO}: define $S$ and $S_{kat}$ (should be fairly similar to the $\Lambda$ flow case)

\subsubsection{$\Lambda$ flow}

From the $T$-$\Lambda$ flow [\ref{PMFRGXYZTLambda}] (see line 1732) we get:
\begin{equation}
  iG{(\omega,T,x,\Lambda)}^{-1}_{a,b} = \frac{\omega^2 \delta_{a,b} + i \omega \Sigma^\Lambda_{ab}(x,\omega) + \Lambda^2}{\omega}
\end{equation}

The function S (the ``single scale propagator'', see [\ref{PFFRGReview}], Eq. 56 and footnote 11)
is
\begin{equation}
  S = -\left. \frac{\partial G}{\partial \Lambda} \right|_{\Sigma^\Lambda = const} ,
\end{equation}
i.e., the derivative of this $G$ with respect to $\Lambda$, keeping $\Sigma$ constant.


So in this case
\begin{equation}
  S_{ab} = \frac{2\Lambda}{\omega} G_{ac} G_{cb}
\end{equation}
Where we left understood the arguments $(\omega,T,x,\Lambda)$ both for functions $S_{ab}$ and $G_{ab}$.

Notice that in the Flow equations $S_{kat}$ is actually used,
which should improve the quality of the truncation (\emph{Katanin truncation}).
In this case, we use the full derivative,
where we also include the derivative $\partial_\Lambda \Sigma^\Lambda$

\begin{equation}
  S_{kat} = \frac{\partial G}{\partial \Lambda} = G_{ac} \left( \frac{2\Lambda}{\omega}\delta_{cd}  +i \partial_\Lambda \Sigma_{cd}^\Lambda(x,\omega)\right)G_{db},
\end{equation}
where for $\partial_\Lambda \Sigma^\Lambda$ we can use the derivative we compute in the flow equations.

\section{Flow Equations}
TODO: this mimics Eq.4 from Noah's notes, and Eq. 72 from Siebe's,
     but
\begin{equation}
  \frac{d}{d\Lambda} \Gamma^{ij}_{f_1,f_2,f_3,_f4}(\omega_1,\omega_2,\omega_3,\omega_4) =
\end{equation}



\section{Majorana Algebra}
From the Majorana sub-package,
using the \textsf{print\_Hamiltonian\_latex()} function
in the \textsf{src/Hamiltonian.jl} file,
and then slightly massaging the output
in a trivial way:
\begin{equation}
  \begin{aligned}
    H = -\frac{1}{4} \sum_{i,j} \eta^x_i \eta^y_i \eta^x_j \eta^y_j  \\
+2 \eta^x_i \eta^y_i \eta^z_j \theta^z_j \\
+\eta^x_i \eta^z_i \eta^x_j \eta^z_j \\
+\sqrt{3} \eta^x_i \eta^z_i \eta^y_j \theta^x_j \\
+\eta^x_i \eta^z_i \eta^y_j \theta^z_j \\
+3 \eta^x_i \theta^x_i \eta^x_j \theta^x_j  \\
- \sqrt{3} \eta^x_i \theta^x_i \eta^x_j \theta^z_j \\
+\sqrt{3} \eta^x_i \theta^x_i \eta^y_j \eta^z_j \\
- \sqrt{3} \eta^x_i \theta^z_i \eta^x_j \theta^x_j \\
+\eta^x_i \theta^z_i \eta^x_j \theta^z_j \\
-\eta^x_i \theta^z_i \eta^y_j \eta^z_j \\
+\sqrt{3} \eta^y_i \eta^z_i \eta^x_j \theta^x_j  \\
- \eta^y_i \eta^z_i \eta^x_j \theta^z_j \\
+\eta^y_i \eta^z_i \eta^y_j \eta^z_j \\
+\sqrt{3} \eta^y_i \theta^x_i \eta^x_j \eta^z_j  \\
+3 \eta^y_i \theta^x_i \eta^y_j \theta^x_j \\
+\sqrt{3} \eta^y_i \theta^x_i \eta^y_j \theta^z_j  \\
+\eta^y_i \theta^z_i \eta^x_j \eta^z_j \\
+\sqrt{3} \eta^y_i \theta^z_i \eta^y_j \theta^x_j \\
+\eta^y_i \theta^z_i \eta^y_j \theta^z_j  \\
+2 \eta^z_i \theta^z_i \eta^x_j \eta^y_j \\
+4 \eta^z_i \theta^z_i \eta^z_j \theta^z_j
  \end{aligned}
\end{equation}

After visual inspection, all terms match with Eq.(13)
in Siebe's notes, except the sign (or order)
of $\eta^y_i \theta^z_i \eta^x_j \eta^z_j$,
which in Siebe's notes is reported as
$\eta^y_i \theta^z_i \eta^z_j \eta^x_j$
instead (still with a plus sign).

\section{Bibliography}


\begin{enumerate}
\item General PFFRG review: \href{https://arxiv.org/pdf/2307.10359}{https://arxiv.org/pdf/2307.10359} \label{PFFRGReview}

\item PMFRG with XYZ anisotropy (code): \href{https://github.com/NoahHassan/pmfrg\_xyz}{https://github.com/NoahHassan/pmfrg\_xyz}
\item PMFRG xyz $T$-$\Lambda$ flow unified code: \href{https://github.com/mmesiti/pmfrg\_xyz/blob/a2a8d10af3b27aeb67f0d2a1b67816cc5d309d96/src/PMFRG\_xyz.jl} {https://github.com/mmesiti/pmfrg\_xyz/blob/a2a8d10af3b27aeb67f0d2a1b67816cc5d309d96/src/PMFRG\_xyz.jl} \label{PMFRGXYZTLambda}
\end{enumerate}

\end{document}
